\documentclass[10pt]{article}
\usepackage{cornell-notes}
\title{Chapter 9: Voltage Regulation and Power Conversion}\author{}\date{}
\setDocTitle{Chapter 9: Voltage Regulation and Power Conversion}\setDocSubtitle{Electronics}\setDocOwner{Electronics Cornell Notes}\setDocDate{\today}
\setCornellCollection{Electronics Cornell Notes}\setCornellUnitType{Chapter}\setCornellUnitNumber{9}\setCornellUnitTitle{Voltage Regulation and Power Conversion}
\hypersetup{pdftitle={Chapter 9: Voltage Regulation and Power Conversion},pdfsubject={Cornell notes},pdfauthor={}}
\begin{document}\maketitle
\begin{cornell}\noterow{Core idea}{Power conversion controls energy flow through switching, feedback, magnetics, and thermal limits.}\noterow{Validation}{Check efficiency, ripple, current limits, stability, and temperature.}\end{cornell}
\begin{summarybox}{Section summary}A converter is reliable only across its electrical and thermal operating range.\end{summarybox}
\end{document}
